\documentclass{article}
\usepackage{amsmath, amssymb, amsthm, graphicx, rotating, float, tikz, listings, xcolor}

\usepackage[a4paper, margin=2cm]{geometry}
\usepackage{amsmath, mathtools, amssymb}
\usepackage{titlesec}
\usepackage[inline]{enumitem}


\title{Cloud and Concurrent Programming | Video Player Plan}
\author{a1967663 | November 2026}
\date{}

\begin{document}
\maketitle

\section*{Captions and Subtitles (Client)}
The core feature of this assignment is about subtitles, so this should be the main priority. We're required to include a feature that translates video to subtitles in real time. This means calling an external api at \textit{https://api.openai.com/v1/audio/transcriptions}, and showing the subtitles using frame processors as seen in Assignment 2.

\subsection*{Displaying captions}
The client grabs a cue from the server at its current time-frame. This cue is processed by the \textit{SubtitleRenderer} class, which generates the text, and how long it is displayed for. Then, a new frame processor \textit{FrameSubtitle} class recieves the rendered cue strip, and burns it into the oncoming frame.

\subsection*{Subtitle styling and user requirements}
To make subtitles visibly engaging and feel natural across multiple backgrounds, we need to style how the frame processor handles drawing subtitles. Two new files \textbf{SubtitleColours.java} and \textbf{SubtitleFont.java} were first initialised using Java's GUI library \textit{awt} to make modifying text easier by hand. Some other optimisation feature includes bit manip on argb payload, and re-using cue strips to avoid unnecessary recomputation
\\\\
To handle the background and how much characters should be displayed, \textbf{SubtitleRenderer.java} draws the cue with the given colour and font, then \textbf{FrameSubtitle} processor unpacks that and burns it into the next frame.

\section*{Video Player Window control (Client)}
The current window of the video player is ever only in full screen or sits awkwardly in the corner of your screen on a monitor, and you cant drag the window anywhere. It is important that the user always has control of where the window should be. Therefore, we made the window draggable, and resizable any time.

\subsection*{Approach to movable window}
My initial approach is to reprocess the \textbf{VideoPlayerApp} class to create a top border on the current window, then create a click and drag detection. As the user clicks, hold and moves while mousedown on the top bar, the window moves relative to the mouse while clicking and dragging the window.

\subsection*{Approach to resizable window}
Similar to the draggable window approach, we add mouse detectors all around the edges and borders of the window to detect if it is about to be interacted with. Then, as the window detects mouse movement while dragging the edge of the borders, it used the \textit{Stage} class and calls its \textit{set} methods

\section*{Advanced Scrubbing (Client)}
I'm taking heavy inspiration from Youtube here, and the main feature I use is their advanced scrubbing system; being able to click on the bar and drag around while seeing the frame for that specific video. I'd love to research a way to implement this soon.

\section*{Advanced Audio control (Client)}
Another nice control to have is to change the volume of the video. Instead of a binary on/off button, I want dynamic audio control.

\section*{Magnifier control (Client)}
A feature I sometimes see myself inadvertently doing when using any video player or photos app is zooming in/out of something small in a video. Perhaps this could be done if I have enough time.

% \section*{Other settings options}
% Other setting and customisation including resizing, moving and changing the opacity of translation. Also, applying and removing filters in real-time instead of having to pre-run a command.

\section*{Captions and Subtitles (Server)}
The subtitles on the server-side will handle OpenAPI's whisper model, as well as converting the response to cues, which will later be grabbed from by the client.

\subsection*{Real-time versus Pre-processing subtitles}
We need to decide whether we should analyse all the sound file at once to grab the subtitles before the video plays, or process each relevant subtitles during playtime. As pre-processing the entire video for subtitles is easier, the user must pay the cost of having to wait \textbf{a long time}. Instead, we'll process as the video gets watched.

\subsection*{Subtitle processing timing}
We can't just process the next couple seconds of where we currently are as that would make each caption generated be relative to the position we played it as, making it difficult to optimise through caching in the future. Instead, we'll naively partition the video into $n$ different chunks, each of $N$ seconds long. As we enter each chunk, the subtitles is generated for that chunk, and cached to the server so it is never called again. 
\\\\
However there is only one complex issue that this causes: If a word were long enough to be between two chunks, then it'll get cut off and look unnatural. Currently I don't know how to fix this, but maybe I can figure something out in the future

\subsection*{}
% \begin{figure}[H]
%     \centering
%     \includegraphics[width=0.5\linewidth]{img1.png}
%     \caption{Very expensive minesweeper, costly interaction with neighbours}
% \end{figure}

\end{document}
